\label{sec:support}
\subsection{The network is a difference of support functions}
A bias-free \textsc{ReLU} network is positively homogeneous: $f(x)=|x|f(\hat x)$. Under $\gamma$ the radius and the direction are
independent, so $\E f=\E|x|\cdot\int_{S^{n-1}}f\,d\sigma$: the problem lives on the sphere. Every positively homogeneous piecewise-linear
function is a difference of support functions of polytopes (the tropical description of Zhang--Naitzat--Lim). Write
$h_K(x)=\max_{y\in K}\langle y,x\rangle$.
\begin{proposition}[tropical decomposition; proved]\label{prop:tropical}
$\relu(h_A-h_B)=h_{\mathrm{conv}(A\cup B)}-h_B$. Hence, inductively, $f=h_{P_+}-h_{P_-}$ for two polytopes $P_\pm$ built from the rows by
Minkowski sums and convex hulls. A first hidden layer contributes zonotopes $\sum_ac_a[0,w_a]$.
\end{proposition}
\begin{theorem}[the mean is a difference of first intrinsic volumes; proved]\label{thm:V1}
\[
\E f=\frac{\E|x|}2\big(w(P_+)-w(P_-)\big)=\frac{V_1(P_+)-V_1(P_-)}{\sqrt{2\pi}},
\]
where $w$ is the mean width and $V_1$ the first intrinsic volume. Moreover $V_1(P)=\sum_{\text{edges }e}|e|\,\gamma(e,P)$, the edge
length times the external angle. Under the normal-fan duality this is stage 11's border formula,
$\E f=(2\pi)^{-1/2}\sum_F\kappa_F\alpha_F$:
\begin{itemize}[leftmargin=*]
\item the cells of the activation tiling are normal cones of vertices;
\item the walls are normal cones of edges;
\item the kink $\kappa_F$ is the edge length $|e|$;
\item $\alpha_F$ is the external angle.
\end{itemize}
\end{theorem}
\begin{proof}
$\E h_K(x)=\E\sup_{y\in K}\langle y,x\rangle=V_1(K)/\sqrt{2\pi}$ (Sudakov--Tsirelson), and polar coordinates give the mean-width form.
The edge formula is standard (Schneider, Section 4.2).
\end{proof}
Checked in $\R^2$ with $8$ hidden units and a \textsc{ReLU} output.
\begin{itemize}[leftmargin=*]
\item The identity of Proposition~\ref{prop:tropical} holds to $6\times10^{-16}$.
\item The quadrature mean $0.16269286$ equals $\frac{\E|x|}2\cdot\frac{\mathrm{per}(P_+)-\mathrm{per}(P_-)}\pi$ to all printed digits.
\end{itemize}
The packing and convex-geometry literature is therefore natively about our objects. A network is a pair of polytopes, its mean is an
intrinsic volume, its walls are a normal fan, and its gates are a hyperplane tessellation of the sphere. For the last, two inputs at
angle $\theta$ are separated by a fraction $\theta/\pi$ of the hyperplanes (stage 11).

\subsection{Each layer evaluates the support function of a zonoid}
For a law $\mu$ on $\R^n$, the zonoid $Z(\mu)=\{\E[Xg(X)]:\ 0\le g\le1\}$ (Koshevoy--Mosler) has support function $h_{Z(\mu)}(w)=\E\langle w,X\rangle_+$.
\begin{proposition}[the layer map is support-function evaluation; proved]
$\E h_{l,a}=h_{Z(\mu_{l-1})}(w_{l,a})$. The next layer's means are the values of one convex body's support function at the $n$ rows.
The Gaussian closure replaces $Z(\mu_{l-1})$ by the zonoid of $N(\mu,C)$, whose support function is $s\,m_1(r)$. For centred laws this
zonoid is the ellipsoid $C^{1/2}B/\sqrt{2\pi}$. In general it is an ``egg'' interpolating between that ellipsoid and the segment to the mean.
\end{proposition}
\begin{proposition}[the curvature of the zonoid is the wall density; proved]\label{prop:curv}
$\nabla_w^2\,h_{Z(\mu)}(w)=\E[\delta(\langle w,X\rangle)\,X\otimes X]$. In the plane, the radius of curvature $h+h''$ is the density of the projected
law at the wall, weighted by the transverse second moment.
\end{proposition}
Checked for a three-component Gaussian mixture in $\R^2$: the finite-difference Hessian and the wall integral agree to six digits.
This joins three pictures:
\begin{itemize}[leftmargin=*]
\item stage 12's wall-density law (the lower tier's activity is the zonoid's curvature);
\item stage 11's boundary layer;
\item the billiard paper's positivity $g+g_{\theta\theta}>0$.
\end{itemize}
In a $2$-plane, the angular harmonics of $h_Z$ away from the Gaussian egg are the cumulant corrections by order: $\kappa_3$ feeds odd
harmonics and $\kappa_4$ even ones. Non-Gaussianity is non-ellipticity of zonoids.

\subsection{Klartag's ellipsoid in our coordinates (dictionary; analogy, labelled)}
\begin{center}\small
\begin{tabular}{>{\raggedright\arraybackslash}p{5.0cm}>{\raggedright\arraybackslash}p{8.8cm}}\toprule
Klartag's process & the network and the chain\\\midrule
ellipsoid $\{A_tx\cdot x<1\}$, $n(n+1)/2$ degrees of freedom & the closure's Gaussian state $(\mu,C)$; its $n(n+1)/2$ covariance degrees of freedom\\
lattice point on the boundary (contact) & a frozen (resolved) gate, $|r_a|\gg1$; the ball no longer crosses that wall\\
free directions $F_t$, $\dim F_t\ge n(n+1)/2-|\text{contacts}|/2$ & unresolved gates, a fraction $\theta/\pi\approx3/\tau$ per layer (stage 12)\\
drift of $\log\det A_t$, $-\dim F_t/2\|A\|^2$ & the information rate of the gate filtration: log-volume of the input cell lost per layer\\
contacts accrue at rate about $e^{n^2t/8}$ & gates freeze like $\sqrt T$, with $T\approx\theta^2/2$ the internal temperature\\
sticky constraint $\langle A_t,x\otimes x\rangle=1$ & an exact invariant the chain must keep: the dilation package, the Euler--Stein ladder\\\bottomrule
\end{tabular}
\end{center}
The structural content is the sticky principle. A process that is only approximately right gains most by holding exactly the
constraints along which it cannot damp its own errors. In Klartag's case these are the contacts, which keep the ellipsoid lattice-free.
In ours they are the critical (charge) directions of Section~\ref{sec:dilation}.

\paragraph{Spherical codes and codegrees (analogy).} The rows of $W_1$ form a random spherical code, and the first layer's cells are
the $2^n$ images of orthants. In the packing work of Campos, Jenssen, Michelen and Sahasrabudhe, the nibble's ``random-like''
heuristic is made exact by controlling \emph{codegrees}. In ours, the closure (independent, dichotomized-Gaussian gates) is the random-like heuristic, and it is
accurate when pair couplings are small. Stage 12 found the bulk couplings $O(n^{-1/2})$ and incoherent. The one coherent,
uncontrolled ``codegree'' is the Mattis (dilation) component. It is again the charge.
